%% 第6巻 巻末（こたえ・記号カード・定数と単位・公式・チェック表・つぎの巻へ・さくいん・奥付）（第14回編集委員会）
%% 第4巻・第5巻の巻末と同じ形。記号カードの一覧・書くときのちゅうい・定数と単位・公式・電卓と表計算と Python・つぎの巻へを、第6巻の本文に合わせて書いた
%% 第6巻はルビなし（\R は使わない）。字の色は bkTheme でなく \bkTT。つぎの巻へのカードは第7巻の色（bkVol7・bkDeep7）
\makeatletter
%% ---------------- こたえ ----------------
\bkAnswerPages

%% ---------------- 記号カード ----------------
%% 前の巻の記号で、この巻でもよく使うもの（章の番号は、はじめて出た巻の章。この巻の章の \bkkigocard は自動でうしろに集まる）
%% 6-1 のカードは記号を \fontsize{14}{16} で大きく見せている。一覧の記号の列（18 mm）には入らないので、
%% 一覧の中だけ、14pt の指定をそのときの字の大きさ（\footnotesize）に読みかえる（章の本文のカードはそのまま）
\let\bk@origfontsize\fontsize
\def\bk@fsfourteen{14}%
\protected\def\fontsize#1#2{\def\bk@fsarg{#1}\ifx\bk@fsarg\bk@fsfourteen\expandafter\bk@origfontsize\expandafter{\f@size}{\f@baselineskip}\else\bk@origfontsize{#1}{#2}\fi}%
\bkKigoPages{%
\bkkigorow{$\Delta$}{デルタ}{変化量。「あと$-$まえ」。この巻の $\Delta t$ は、オイラー法の1歩の短い時間}{2-3}%
\bkkigorow{$D$}{ディー}{判別式 $b^2-4ac$。この巻では、特性方程式の解が2つ・重解・複素数のどれかを見分ける}{3-10}%
\bkkigorow{$\theta$}{シータ}{角の大きさ。微分・積分や $e^{j\theta}$ ではラジアン}{4-2}%
\bkkigorow{$\tan^{-1}$}{アークタンジェント}{$\tan$ がその数になる角。右半分（$-90^\circ$〜$90^\circ$）の角しか答えない}{4-6}%
\bkkigorow{$\omega$}{オメガ}{角速度。この巻では、ゆれの角振動数（rad/s）}{4-8}%
\bkkigorow{$\dfrac{\dd y}{\dd x}$}{ディーワイ・ディーエックス}{$y$ を $x$ で微分したもの。$\dfrac{\dd T}{\dd t}$ は「$T$ の変わる速さ」}{5-3}%
\bkkigorow{$\dot{x}$、$\ddot{x}$}{エックス・ドット、エックス・ツードット}{時間で1回微分（速度）・2回微分（加速度）。$\dot x=\dfrac{\dd x}{\dd t}$}{5-5}%
\bkkigorow{$\approx$}{ニアリーイコール}{ほぼ等しい（近似）}{5-8}%
\bkkigorow{$\displaystyle\int_a^b$}{インテグラル}{$a$ から $b$ まで、細かく分けてたす（定積分）}{5-10}%
\bkkigorow{$e$}{イー（ネイピア数）}{$2.718\cdots$。$(e^x)'=e^x$}{5-14}%
\bkkigorow{$\ln x$}{エル・エヌ・エックス}{自然対数。$e$ を何乗すると $x$ になるか}{5-14}%
\bkkigorow{$\phi$}{ファイ}{初期位相。この巻では、位相のおくれ（6-13）にも使う}{5-15}%
\bkkigorow{$n!$}{エヌのかいじょう}{1から $n$ までの整数をすべてかけた数。$0!=1$}{5-17}%
}
\let\fontsize\bk@origfontsize

%% 記号カードのつづき：ギリシャ文字（この巻での意味に合わせた）
\def\bkbkhead#1{\par\addvspace{4mm}\noindent{\sffamily\bfseries\large\color{\bkTT}#1}\par\nobreak\vspace{2mm}}
\par\noindent\begin{minipage}{\linewidth}
\bkbkhead{ギリシャ文字の読み方}
{\small この巻では、時定数 $\tau$、減衰比 $\zeta$、固有角振動数 $\omega_{\mathrm n}$ など、ギリシャ文字がたくさん出てきます。大文字・小文字・読み方と、この巻での意味です（かっこの中は章）。}\par\vspace{2mm}
\begin{bkcenter}
\footnotesize
\begin{tblr}{colspec={Q[c,m,wd=7mm] Q[c,m,wd=7mm] Q[l,m,wd=17mm] Q[l,m,wd=30mm] Q[c,m,wd=7mm] Q[c,m,wd=7mm] Q[l,m,wd=17mm] Q[l,m,wd=30mm]}, row{1}={bg=bkVol!14, font=\sffamily\bfseries\footnotesize}, hlines={0.4pt,bkLine}, vline{5}={0.8pt,bkVol}, rowsep=0.8pt, colsep=2.5pt}
大 & 小 & 読み方 & この巻での意味 & 大 & 小 & 読み方 & この巻での意味 \\
$\mathrm{A}$ & $\alpha$ & アルファ & へる速さ $\zeta\omega_{\mathrm n}$（6-7）、方程式の解（6-8） & $\mathrm{N}$ & $\nu$ & ニュー & 振動数 \\
$\mathrm{B}$ & $\beta$ & ベータ & 方程式の解（6-8） & $\Xi$ & $\xi$ & クシー & （変数） \\
$\Gamma$ & $\gamma$ & ガンマ & （係数） & $\mathrm{O}$ & $o$ & オミクロン & ― \\
$\Delta$ & $\delta$ & デルタ & $\Delta t$ は短い時間（6-2）。$\delta$ は対数減衰率（6-7） & $\Pi$ & $\pi$ & パイ & 円周率 \\
$\mathrm{E}$ & $\varepsilon$ & イプシロン & 小さな数（6-12） & $\mathrm{P}$ & $\rho$ & ロー & 密度（6-3） \\
$\mathrm{Z}$ & $\zeta$ & ゼータ & 減衰比（6-7） & $\Sigma$ & $\sigma$ & シグマ & $e^{(\sigma+j\omega)t}$ の $\sigma$（6-10） \\
$\mathrm{H}$ & $\eta$ & イータ & 粘度（6-5） & $\mathrm{T}$ & $\tau$ & タウ & 時定数（6-4） \\
$\Theta$ & $\theta$ & シータ & 振り子のふれ（6-6）、偏角（6-9） & $\Upsilon$ & $\upsilon$ & ウプシロン & ― \\
$\mathrm{I}$ & $\iota$ & イオタ & ― & $\Phi$ & $\phi$ & ファイ & 初期位相（6-6）、位相のおくれ（6-13） \\
$\mathrm{K}$ & $\kappa$ & カッパ & 弦や柱の形の数（6-6） & $\mathrm{X}$ & $\chi$ & カイ & ― \\
$\Lambda$ & $\lambda$ & ラムダ & 1秒あたりの平均の回数（6-4） & $\Psi$ & $\psi$ & プサイ & （角） \\
$\mathrm{M}$ & $\mu$ & ミュー & マイクロ（\unit{\micro F}） & $\Omega$ & $\omega$ & オメガ & 角振動数（6-6）。$\Omega$ は抵抗の単位 \\
\end{tblr}
\end{bkcenter}
{\footnotesize 大文字が英語のアルファベットと同じ形の文字（A、B、E など）は、記号としてはほとんど使いません。$\phi$ は $\varphi$ と書くこともあります（4-15 の球座標の角は $\varphi$）。数学の本では、特性方程式の文字に $\lambda$ を使うことがあります（この巻は $s$。6-12）。}
\end{minipage}

%% 記号カードのつづき：手で書くとき・読むときのちゅうい（各章の「まちがえやすいところ」から）
\Needspace{12\baselineskip}
\bkbkhead{記号を書くときのちゅうい}
{\small 手で書くと、ほかの字や記号とまぎれやすいものと、読みまちがえやすい書き方です。}\par\vspace{2mm}
{\small
\begin{longtblr}{colspec={Q[c,m,wd=19mm] Q[l,m,wd=22mm] Q[l,m,wd=85mm]}, rowhead=1, row{1}={bg=bkVol!14, font=\sffamily\bfseries\footnotesize}, hlines={0.4pt,bkLine}, rowsep=0.8pt}
記号 & まぎれやすい & 書くとき・読むときのコツ \\
$\dfrac{\dd T}{\dd t}$ & $\dfrac{T}{t}$ & 「$T$ の変わる速さ」。$\dd$ で約分しない（6-1） \\
$\dot x$、$\ddot x$ & $x$、$\bar x$ & 点は文字の真上に。2回微分は点2つ（6-6） \\
$y_n$\ \ $x_1$ & ― & $y_n$ は $n$ 歩め（6-2）、$x_1$ は1番めの物（6-14） \\
$\omega_{\mathrm n}$ & $\omega_n$ & n は名前（natural）の頭文字で、立てた字（6-6） \\
$e^{-t/\tau}$ & $e^{-\tau t}$ & $\tau$ でわる。$\tau=\qty{2}{s}$ なら $e^{-t/2}$（6-4） \\
$c$ & $c$ & 6-5 は2乗の抵抗の係数〔N·s$^2$/m$^2$〕、6-7 からはダンパーの減衰係数〔N·s/m〕（6-7） \\
$\zeta$ & 2、$z$ & 上の横線から左下へ流し、下で右へはねる（6-7） \\
$j$ & $i$ & 虚数単位。$i$ は電流に使うので $j$。高校では $i$（6-8） \\
$3+4j$ & ― & 虚部は4（$4j$ ではない）。$j^2$ はすぐ $-1$ に（6-8） \\
$\bar z$ & $-z$ & 上の横線は共役。$a+bj$ なら $a-bj$（6-8） \\
$|z|$ & ― & $\sqrt{a^2+b^2}$。$\sqrt{\phantom{2}}$ の中に $j$ を入れない（6-9） \\
$\arg z$ & $\tan^{-1}\dfrac{b}{a}$ & $-180^\circ<\theta\leqq180^\circ$。点の区画を見る（6-9） \\
$e^{j\theta}$ & $e^{j60}$ & 肩の $\theta$ はラジアン。60°なら $e^{j\pi/3}$（6-10） \\
$\cosh x$ & $\cos x$ & $(\cosh x)'=\sinh x$（$-$ はつかない）。$\cosh^2x-\sinh^2x=1$（和ではなく差）（6-11） \\
$\mathrm{Re}[\ ]$ & ― & 実部をとる。答えに $j$ を残さない（6-12・6-13） \\
$\leftrightarrow$ & $=$ & 「同じ役わり」の意味。等しいのではない（6-15） \\
$V$〔V〕 & $V$\,V（VV） & 記号と単位が並ぶときは、単位を〔\ 〕でくくる。「電圧 $V$〔V〕」「時定数 $\tau$〔s〕」（6-0） \\
\end{longtblr}\par}

%% ---------------- 定数・単位・公式 ----------------
\bkbacktitle{よく使う定数・単位・公式}
\bksubsection{1}{この巻で使う単位}
{\small
\begin{longtblr}{colspec={Q[l,m,wd=24mm] Q[l,m,wd=34mm] Q[l,m,wd=52mm] Q[c,m,wd=12mm]}, rowhead=1, row{1}={bg=bkVol!14, font=\sffamily\bfseries\footnotesize}, hlines={0.4pt,bkLine}, rowsep=1.4pt}
量 & 単位（読み方） & ほかの単位との関係 & 章 \\
時間 & s、min、h & $\qty{60}{s/min}$、$\qty{3600}{s/h}$。\ \qty{1}{ms}は $\frac{1}{1000}$\,s & 6-3 \\
温度 & ℃、K（ケルビン） & 温度差は \qty{1}{K}＝1\,℃ と同じ大きさ & 6-3 \\
電圧・電流 & V（ボルト）、A（アンペア） & $\qty{1}{mA}=\qty{0.001}{A}$ & 6-1 \\
電荷 & C（クーロン） & $\qty{1}{C}=\qty{1}{A.s}$（電流×時間） & 6-2 \\
抵抗 & $\Omega$（オーム） & $\qty{1}{\ohm}=\qty{1}{V/A}$、\ $\qty{1}{k\ohm}=\qty{1000}{\ohm}$ & 6-2 \\
静電容量 & F（ファラド） & $\qty{1}{F}=\qty{1}{C/V}$、\ $\qty{1}{\micro F}=\qty{0.000001}{F}$ & {6-2\\6-4} \\
インダクタンス & H（ヘンリー） & $\qty{1}{H}=\qty{1}{V.s/A}$、\ $\qty{1}{mH}=\qty{0.001}{H}$ & 6-4 \\
時定数 & s & $\Omega\times\mathrm F=\mathrm s$（$RC$）、\ $\dfrac{\mathrm H}{\Omega}=\mathrm s$（$\dfrac{L}{R}$） & {6-2\\6-4} \\
熱容量 & J/K & 温度を\qty{1}{K}上げるのに要る熱 & {6-3\\6-4} \\
熱抵抗 & K/W & \qty{1}{W}の熱を流すのに要る温度差 & 6-4 \\
電力・熱流 & W（ワット） & $\qty{1}{W}=\qty{1}{J/s}$ & 6-4 \\
圧力 & Pa（パスカル） & $\qty{1}{Pa}=\qty{1}{N/m^2}$、\ $\qty{1}{hPa}=\qty{100}{Pa}$ & 6-3 \\
ばね定数 & N/m & \qty{1}{m}のばすのに要る力 & 6-6 \\
抵抗の係数 & N·s/m（$b$）、\ N·s$^2$/m$^2$（$c$） & 力 $bv$、$cv^2$ が N になる & 6-5 \\
減衰係数 & N·s/m & ダンパーの力 $c\dot x$ が N になる & 6-7 \\
角振動数 & rad/s & $\omega=2\pi f$ & 6-6 \\
振動数 & Hz（ヘルツ） & 1秒に何回。$f=\dfrac{1}{T}$、\ $\qty{1}{kHz}=\qty{1000}{Hz}$ & {6-6\\6-13} \\
\end{longtblr}\par}

\bksubsection{2}{よく使う数}
\begin{bkcenter}
\small
\begin{tblr}{colspec={Q[l,m,wd=46mm] Q[l,m,wd=60mm] Q[c,m,wd=18mm]}, row{1}={bg=bkVol!14, font=\sffamily\bfseries\footnotesize}, hlines={0.4pt,bkLine}, rowsep=1pt}
名前 & 値 & 章 \\
ネイピア数 $e$ & 2.71828… & 5-14 \\
$e^{-1}$ & 0.368（時定数 $\tau$ で残る差 36.8\,\%） & 6-4 \\
$\ln2$ & 0.693（半減期 $=\tau\ln2$） & 6-4 \\
円周率 $\pi$、\ $2\pi$ & 3.14159…、\ 6.28318… & 6-6 \\
$\dfrac{1}{\sqrt2}$ & 0.707 & {6-12\\6-13} \\
重力加速度 $g$ & \qty{9.8}{m/s^2} & 6-5 \\
水の比熱 & 1\,gあたり約\qty{4.2}{J/K} & 6-4 \\
水の密度 & \qty{1000}{kg/m^3} & 6-3 \\
大気の圧力（海面） & \num{101300}\,Pa（1013\,hPa） & 6-3 \\
\end{tblr}
\end{bkcenter}

\bkbkhead{この巻でくり返し使う場面と数}
{\small 同じ場面は、どの章でも同じ数で計算しています。}\par\vspace{1mm}
{\small
\begin{longtblr}{colspec={Q[l,m,wd=24mm] Q[l,m,wd=86mm] Q[c,m,wd=14mm]}, rowhead=1, row{1}={bg=bkVol!14, font=\sffamily\bfseries\footnotesize}, hlines={0.4pt,bkLine}, rowsep=1.6pt}
場面 & 数 & 章 \\
冷めるお湯 & 部屋20\,℃、はじめ80\,℃、$\dfrac{\dd T}{\dd t}=-0.05\,(T-20)$（$t$ は分）。$T=20+60e^{-0.05t}$、$\tau=20$ 分。10分後 56.39\,℃（オイラー法 $\Delta t=1$ 分では 55.92\,℃）。50\,℃まで約13.9分 & {6-1\\6-2\\6-4} \\
RC 回路 & $E=\qty{10}{V}$、$R=\qty{20}{k\ohm}$、$C=\qty{100}{\micro F}$、$RC=\qty{2}{s}$。充電 $V=10\,(1-e^{-t/2})$ & {6-2\\6-4} \\
ばね & $m=\qty{0.5}{kg}$、$k=\qty{8}{N/m}$、$\omega_{\mathrm n}=4$\,rad/s、周期 約\qty{1.57}{s}。臨界 $c_{\mathrm c}=\qty{4}{N.s/m}$。$c=\qty{0.4}{N.s/m}$ で $\zeta=0.1$、$\omega_{\mathrm d}=3.980$\,rad/s、対数減衰率 0.631 & {6-6\\6-7\\6-12\\6-13} \\
RLC 回路 & $L=\qty{10}{mH}$、$C=\qty{1}{\micro F}$、$\omega_{\mathrm n}=\num{10000}$\,rad/s（約\qty{1592}{Hz}）。臨界 $R_{\mathrm c}=\qty{200}{\ohm}$。$R=\qty{20}{\ohm}$ で $\zeta=0.1$ & {6-6\\6-7\\6-12\\6-13} \\
スカイダイバー & $m=\qty{80}{kg}$、$c=\qty{0.26}{N.s^2/m^2}$、終端速度 $v_{\mathrm t}\approx\qty{54.9}{m/s}$ & {6-5\\6-11} \\
3-10 の2つの式 & $s^2+3s+2=0$ の解は $s=-1$、$-2$（ゆれずにもどる）。$s^2+2s+5=0$ の解は $s=-1\pm2j$（ゆれながらもどる） & {6-8\\6-12} \\
2つの水そう & $\dfrac{\dd h_1}{\dd t}=-0.1\,(h_1-h_2)$、$\dfrac{\dd h_2}{\dd t}=0.1\,(h_1-h_2)$。$h_1=10+10e^{-0.2t}$ & 6-14 \\
2つの台車 & $m=\qty{1}{kg}$、$k=\qty{100}{N/m}$、$K=\qty{10}{N/m}$。$\omega_1=10$、$\omega_2\approx10.95$\,rad/s（うなり） & 6-14 \\
\end{longtblr}\par}

\bksubsection{3}{1回微分の式の公式}
{\small
\begin{longtblr}{colspec={Q[l,m,wd=28mm] Q[l,m,wd=84mm] Q[c,m,wd=12mm]}, rowhead=1, row{1}={bg=bkVol!14, font=\sffamily\bfseries\footnotesize}, hlines={0.4pt,bkLine}, rowsep=1.6pt}
公式・きまり & 式 & 章 \\
微分方程式と解 & 変わる速さが入った方程式。解は関数。代入して、左辺＝右辺を確かめる。初期条件は階数の数だけ要る & 6-1 \\
オイラー法 & $y_{n+1}=y_n+f(y_n)\,\Delta t$（次の値＝今の値＋今の変わる速さ×$\Delta t$）。$\Delta t$ を半分にすると、ずれもほぼ半分 & 6-2 \\
速度を先に進めるオイラー法 & ① $a$ を今の $x$ から出す\quad ② 新しい $v$＝$v+a\,\Delta t$\quad ③ 新しい $x$＝$x+$（新しい $v$）$\times\Delta t$ & 6-2 \\
積分するだけで解ける形 & $\dfrac{\dd y}{\dd t}=f(t)$ なら $y=\displaystyle\int f(t)\,\dd t+C$。\ $\dfrac{\dd^2T}{\dd x^2}=0$ なら $T=C_1x+C_2$ & 6-3 \\
水の深さと圧力 & $\dfrac{\dd p}{\dd z}=-\rho g$（$z$ は上向き）。水では $p=\num{101300}-9800z$〔Pa〕 & 6-3 \\
変数分離 & $\dfrac{\dd u}{\dd t}=-\dfrac{u}{\tau}$ → $\displaystyle\int\frac{\dd u}{u}=\int\Bigl(-\frac{1}{\tau}\Bigr)\dd t$ → $u=Ae^{-t/\tau}$ & 6-4 \\
目標へ近づく形 & $\dfrac{\dd y}{\dd t}=\dfrac{y_\infty-y}{\tau}$ → $y=y_\infty+(y_0-y_\infty)\,e^{-t/\tau}$ & 6-4 \\
半減期 & 差が半分になる時間 $t_{\mathrm h}=\tau\ln2\approx0.693\,\tau$ & 6-4 \\
熱の式 & $C_{\mathrm{th}}\dfrac{\dd T}{\dd t}=P-\dfrac{T-T_{\mathrm a}}{R_{\mathrm{th}}}$。時定数 $\tau=C_{\mathrm{th}}R_{\mathrm{th}}$ & 6-4 \\
コンデンサー & 放電 $\dfrac{\dd V}{\dd t}=-\dfrac{V}{RC}$、\ 充電 $\dfrac{\dd V}{\dd t}=\dfrac{E-V}{RC}$。時定数 $\tau=RC$ & {6-1\\6-2\\6-4} \\
コイル（RL） & $L\dfrac{\dd i}{\dd t}=E-Ri$。最後の電流 $\dfrac{E}{R}$、時定数 $\tau=\dfrac{L}{R}$ & {6-4\\6-5} \\
落ちる物 & $m\dfrac{\dd v}{\dd t}=mg-f$（下向きが $+$）。$f=bv$ なら $v=v_{\mathrm f}\,(1-e^{-t/\tau})$、$v_{\mathrm f}=\dfrac{mg}{b}$、$\tau=\dfrac{m}{b}$ & 6-5 \\
2乗の抵抗 & $f=cv^2$ なら $v_{\mathrm t}=\sqrt{\dfrac{mg}{c}}$、\ $v=v_{\mathrm t}\tanh\dfrac{gt}{v_{\mathrm t}}$ & {6-5\\6-11} \\
体重計 & $N=m\,(g+a)$（上向きが $+$） & 6-5 \\
\end{longtblr}\par}

\vspace{2mm}
\noindent\begin{minipage}{\linewidth}
\noindent{\sffamily\bfseries\small\color{\bkTT}時定数と、残りの差}（6-4）\par\vspace{1mm}
\begin{bkcenter}
\begin{tblr}{colspec={Q[l,m,wd=40mm] *{5}{Q[c,m,wd=15mm]}}, row{1}={bg=bkVol!14, font=\sffamily\bfseries\footnotesize}, hlines={0.4pt,bkLine}, rowsep=1.4pt, cells={font=\small}}
たった時間 & $\tau$ & $2\tau$ & $3\tau$ & $4\tau$ & $5\tau$ \\
残りの差 $e^{-t/\tau}$ & 36.8\,\% & 13.5\,\% & 5.0\,\% & 1.8\,\% & 0.7\,\% \\
目標までに進んだ割合 & 63.2\,\% & 86.5\,\% & 95.0\,\% & 98.2\,\% & 99.3\,\% \\
\end{tblr}
\end{bkcenter}
{\footnotesize 目標にぴったり着くまでには、数学の上では限りなく長い時間がかかります。工学では「$5\tau$ たてば、ほぼ終わり」とみなします。}
\end{minipage}

\bksubsection{4}{2回微分の式（ゆれ）の公式}
{\small
\begin{longtblr}{colspec={Q[l,m,wd=28mm] Q[l,m,wd=84mm] Q[c,m,wd=12mm]}, rowhead=1, row{1}={bg=bkVol!14, font=\sffamily\bfseries\footnotesize}, hlines={0.4pt,bkLine}, rowsep=1.6pt}
公式・きまり & 式 & 章 \\
ばね & $m\ddot x=-kx$、\ $\omega_{\mathrm n}=\sqrt{\dfrac{k}{m}}$、\ 周期 $T=2\pi\sqrt{\dfrac{m}{k}}$ & 6-6 \\
振り子 & $\ddot\theta\approx-\dfrac{g}{L}\theta$、\ $T=2\pi\sqrt{\dfrac{L}{g}}$（ふれが小さいとき） & 6-6 \\
LC 回路 & $L\ddot q=-\dfrac{1}{C}q$、\ $\omega_{\mathrm n}=\dfrac{1}{\sqrt{LC}}$、\ $T=2\pi\sqrt{LC}$ & 6-6 \\
振幅と初期位相 & $x=X\cos(\omega_{\mathrm n}t+\phi)$、\ $X=\sqrt{x_0^{\,2}+\Bigl(\dfrac{v_0}{\omega_{\mathrm n}}\Bigr)^2}$。$\phi$ は点の向きを見て決める & 6-6 \\
つるしたばね & ゆれの中心は、つり合う位置（自然の長さより $\dfrac{mg}{k}$ だけ下） & 6-6 \\
エネルギー & $\dfrac{1}{2}mv^2+\dfrac{1}{2}kx^2$ は一定。$v_{\max}=X\omega_{\mathrm n}$ & {6-6\\6-15} \\
減衰振動の式 & $m\ddot x+c\dot x+kx=0$（ダンパーの力は $-c\dot x$） & 6-7 \\
減衰比 & $\zeta=\dfrac{c}{2\sqrt{mk}}$、\ 臨界 $c_{\mathrm c}=2\sqrt{mk}$。\ $\ddot x+2\zeta\omega_{\mathrm n}\dot x+\omega_{\mathrm n}^2x=0$ & {6-7\\6-12} \\
山のへり方 & つつみ $e^{-\alpha t}$、$\alpha=\dfrac{c}{2m}=\zeta\omega_{\mathrm n}$。対数減衰率 $\delta=\ln\dfrac{\text{前の山}}{\text{あとの山}}\approx2\pi\zeta$ & 6-7 \\
RLC 回路 & $L\ddot q+R\dot q+\dfrac{1}{C}q=0$、\ $\zeta=\dfrac{R}{2}\sqrt{\dfrac{C}{L}}$、\ 臨界 $R_{\mathrm c}=2\sqrt{\dfrac{L}{C}}$ & {6-7\\6-12} \\
特性方程式 & $x=e^{st}$ と置くと、微分は「$s$ をかける」。$ms^2+cs+k=0$ & 6-12 \\
解の形 & $s=-\zeta\omega_{\mathrm n}\pm\omega_{\mathrm n}\sqrt{\zeta^2-1}$。$\zeta<1$ なら $s=-\alpha\pm j\omega_{\mathrm d}$、\ $\omega_{\mathrm d}=\omega_{\mathrm n}\sqrt{1-\zeta^2}$ & 6-12 \\
$A$、$B$ の決め方 & はじめの位置と速度。速度は積の微分で：$x=e^{-t}(A\cos2t+B\sin2t)$ なら $\dot x(0)=-A+2B$ & 6-12 \\
1階の式も & $RC\dot V+V=0$ に $V=e^{st}$ を入れると $RCs+1=0$、\ $s=-\dfrac{1}{RC}=-\dfrac{1}{\tau}$ & 6-12 \\
強制振動 & $m\ddot x+c\dot x+kx=F_0\cos\omega t$。答え＝ゆすりに合うゆれ $x_{\mathrm p}$＋自分のゆれ $x_{\mathrm h}$（やがて消える） & 6-13 \\
置き方 & 右辺が $\cos\omega t$ なら $x_{\mathrm p}=P\cos\omega t+Q\sin\omega t$。$A$、$B$ は $x_{\mathrm h}$ を足してから、はじめの状態で決める & 6-13 \\
複素振幅 & $x=\mathrm{Re}[Xe^{j\omega t}]$、\ $X=\dfrac{F_0}{k-m\omega^2+jc\omega}$、\ 振幅 $|X|=\dfrac{F_0}{\sqrt{(k-m\omega^2)^2+(c\omega)^2}}$ & 6-13 \\
共振 & 振幅の山は $\omega=\omega_{\mathrm n}\sqrt{1-2\zeta^2}$。位相のおくれ $\tan\phi=\dfrac{c\omega}{k-m\omega^2}$（$\omega=\omega_{\mathrm n}$ で90°） & 6-13 \\
Q 値 & $Q=\dfrac{\omega_{\mathrm n}}{\text{半値幅}}$、\ $\zeta$ が小さいとき $Q\approx\dfrac{1}{2\zeta}$ & 6-13 \\
ダンパーなしの共振 & $\ddot x+4x=\cos2t$ の答えの1つは $x=\dfrac{1}{4}t\sin2t$（ゆれがふえ続ける） & 6-13 \\
RLC の共振 & 共振周波数 $f=\dfrac{1}{2\pi\sqrt{LC}}$ & 6-13 \\
2階を1階2つに & $\dot x=v$、\ $\dot v=-\dfrac{k}{m}x$ のように、位置と速度の2つの式にする & 6-14 \\
振動モード & 同じ $m$ の台車2つ・3本のばね：いっしょ $\omega_1=\sqrt{\dfrac{k}{m}}$、反対 $\omega_2=\sqrt{\dfrac{k+2K}{m}}$。どんな動きもモードのたし算 & 6-14 \\
\end{longtblr}\par}

\vspace{2mm}
\noindent\begin{minipage}{\linewidth}
\noindent{\sffamily\bfseries\small\color{\bkTT}特性方程式 $ms^2+cs+k=0$ の3つの場合}（6-12）\par\vspace{1mm}
\begin{bkcenter}
\begin{tblr}{colspec={Q[c,m,wd=14mm] Q[c,m,wd=26mm] Q[c,m,wd=54mm] Q[c,m,wd=28mm]}, row{1}={bg=bkVol!14, font=\sffamily\bfseries\footnotesize}, hlines={0.4pt,bkLine}, rowsep=1.4pt, cells={font=\small}}
判別式 & $\zeta$ と $c$ & 解と、答えの形 & 動き方 \\
$D>0$ & {$\zeta>1$\\$c>2\sqrt{mk}$} & {負の解が2つ\\$Ae^{s_1t}+Be^{s_2t}$} & {ゆっくりもどる\\（過減衰）} \\
$D=0$ & {$\zeta=1$\\$c=2\sqrt{mk}$} & {負の重解\\$(A+Bt)\,e^{st}$} & {ちょうど\\（臨界減衰）} \\
$D<0$ & {$\zeta<1$\\$c<2\sqrt{mk}$} & {$s=-\alpha\pm j\omega_{\mathrm d}$\\$e^{-\alpha t}(A\cos\omega_{\mathrm d}t+B\sin\omega_{\mathrm d}t)$} & {ゆれながら止まる\\（減衰振動）} \\
\end{tblr}
\end{bkcenter}
{\footnotesize 電気の本では、過制動・臨界制動・不足制動ともいいます。行きすぎずにいちばん早く落ち着くのは $\zeta=1$、少し行きすぎてよければ $\zeta=0.7$ くらいがもっと早い（6-12）。}
\end{minipage}

\bksubsection{5}{複素数の公式}
\noindent\begin{minipage}{\linewidth}
\noindent{\sffamily\bfseries\small\color{\bkTT}偏角の求め方}（6-9。$-180^\circ<\theta\leqq180^\circ$）\par\vspace{1mm}
\begin{bkcenter}
\begin{tblr}{colspec={Q[l,m,wd=40mm] Q[l,m,wd=60mm] Q[l,m,wd=24mm]}, row{1}={bg=bkVol!14, font=\sffamily\bfseries\footnotesize}, hlines={0.4pt,bkLine}, rowsep=1.4pt, cells={font=\small}}
点のある所 & 偏角の求め方 & 例 \\
右半分（$a>0$） & $\tan^{-1}\dfrac{b}{a}$ のまま & $1+j$：45° \\
左上（$a<0$、$b\geqq0$） & $\tan^{-1}\dfrac{b}{a}$ に180°をたす & $-1+\sqrt3\,j$：120° \\
左下（$a<0$、$b<0$） & $\tan^{-1}\dfrac{b}{a}$ から180°をひく & $-1-j$：$-135^\circ$ \\
虚軸の上（$a=0$） & $b>0$ なら90°、$b<0$ なら $-90^\circ$ & $-2j$：$-90^\circ$ \\
\end{tblr}
\end{bkcenter}
\end{minipage}\par\vspace{3mm}
{\small
\begin{longtblr}{colspec={Q[l,m,wd=28mm] Q[l,m,wd=84mm] Q[c,m,wd=12mm]}, rowhead=1, row{1}={bg=bkVol!14, font=\sffamily\bfseries\footnotesize}, hlines={0.4pt,bkLine}, rowsep=1.6pt}
公式・きまり & 式 & 章 \\
虚数単位 & $j^2=-1$、\ $j^3=-j$、\ $j^4=1$。$\times j$ は左回りに90°回すこと & 6-8 \\
複素数 & $a+bj$（$a$ は実部 $\mathrm{Re}\,z$、$b$ は虚部 $\mathrm{Im}\,z$）。平面の点 $(a,\ b)$ & {6-8\\6-9} \\
たし算・ひき算 & 実部どうし、虚部どうし。$(3+4j)+(1-2j)=4+2j$（矢印をつなぐ） & {6-8\\6-9} \\
かけ算 & 展開して $j^2=-1$。$(a+bj)(c+dj)=(ac-bd)+(ad+bc)j$ & 6-8 \\
共役 & $\bar z=a-bj$。$(a+bj)(a-bj)=a^2+b^2$ & 6-8 \\
わり算 & 分母と分子に、分母と共役な数をかける。$\dfrac{5}{1+2j}=\dfrac{5(1-2j)}{(1+2j)(1-2j)}=1-2j$。$\dfrac{1}{j}=-j$ & 6-8 \\
負の数の平方根 & $a>0$ のとき $\sqrt{-a}=\sqrt a\,j$。$x^2=-4$ の解は $x=\pm2j$ & 6-8 \\
虚数の解 & $D<0$ なら、解は共役な組 $\dfrac{-b\pm\sqrt{-D}\,j}{2a}$ & 6-8 \\
絶対値 & $|a+bj|=\sqrt{a^2+b^2}$、\ $|z|^2=z\bar z$ & 6-9 \\
極形式 & $z=r(\cos\theta+j\sin\theta)=re^{j\theta}$（$r=|z|$、$\theta=\arg z$） & {6-9\\6-10} \\
かけ算・わり算 & $|zw|=|z|\,|w|$、\ $\arg(zw)=\arg z+\arg w$。わり算は、大きさはわり算、角はひき算 & 6-9 \\
ド・モアブルの定理 & $(\cos\theta+j\sin\theta)^n=\cos n\theta+j\sin n\theta$。\ $(re^{j\theta})^n=r^ne^{jn\theta}$ & {6-9\\6-10} \\
点を回す & 点 $(x,\ y)$ を $\theta$ 回すと $(x\cos\theta-y\sin\theta,\ \ x\sin\theta+y\cos\theta)$ & 6-9 \\
オイラーの公式 & $e^{j\theta}=\cos\theta+j\sin\theta$（単位円の上の角 $\theta$ の点）。$|e^{j\theta}|=1$ & 6-10 \\
特別な角 & $e^{j\pi/2}=j$、\ $e^{j\pi}=-1$（$e^{j\pi}+1=0$）、\ $e^{j2\pi}=1$ & 6-10 \\
肩のたし算 & $e^{ja}e^{jb}=e^{j(a+b)}$ & 6-10 \\
$\cos$ と $\sin$ & $\cos\theta=\dfrac{e^{j\theta}+e^{-j\theta}}{2}$、\ $\sin\theta=\dfrac{e^{j\theta}-e^{-j\theta}}{2j}$ & 6-10 \\
微分 & $\dfrac{\dd}{\dd t}e^{j\omega t}=j\omega\,e^{j\omega t}$（$\omega$ 倍して90°進める） & 6-10 \\
フェーザ & $V_0\sin(\omega t+\phi)$ の矢印は $V_0e^{j\phi}$。同じ $\omega$ の波のたし算は、矢印のたし算 & 6-10 \\
らせん & $e^{(\sigma+j\omega)t}=e^{\sigma t}e^{j\omega t}$。実部は $e^{\sigma t}\cos\omega t$（$\sigma<0$ なら、ゆれながら小さくなる） & 6-10 \\
\end{longtblr}\par}



\bksubsection{6}{双曲線関数の公式}
{\small
\begin{longtblr}{colspec={Q[l,m,wd=28mm] Q[l,m,wd=84mm] Q[c,m,wd=12mm]}, rowhead=1, row{1}={bg=bkVol!14, font=\sffamily\bfseries\footnotesize}, hlines={0.4pt,bkLine}, rowsep=1.6pt}
公式・きまり & 式 & 章 \\
定義 & $\cosh x=\dfrac{e^{x}+e^{-x}}{2}$、\ $\sinh x=\dfrac{e^{x}-e^{-x}}{2}$、\ $\tanh x=\dfrac{\sinh x}{\cosh x}$ & 6-11 \\
2乗の差 & $\cosh^2x-\sinh^2x=1$（点 $(\cosh t,\ \sinh t)$ は双曲線 $x^2-y^2=1$ の上） & 6-11 \\
微分 & $(\cosh x)'=\sinh x$、\ $(\sinh x)'=\cosh x$、\ $(\tanh x)'=1-\tanh^2x=\dfrac{1}{\cosh^2x}$ & 6-11 \\
$\tanh$ のはんい & いつも $-1$ と1のあいだ。くり返さない & 6-11 \\
$y''=a^2y$ の答え & $y=C\cosh at+D\sinh at$（$C=y(0)$、$D=\dfrac{y'(0)}{a}$） & 6-11 \\
懸垂線 & $y=a\Bigl(\cosh\dfrac{x}{a}-1\Bigr)$。たるみが小さいと、放物線 $\dfrac{x^2}{2a}$ に近い & 6-11 \\
\end{longtblr}\par}

\Needspace{9\baselineskip}
\bksubsection{7}{相似の表：ばね・電気・熱}
{\small ちがう分野でも、式が同じ形になります（6-15。物理シリーズ『力学』の表2と同じ組）。}\par\vspace{1mm}
\begin{bkcenter}
\small
\begin{tblr}{colspec={Q[l,m,wd=22mm] Q[c,m,wd=26mm] Q[c,m,wd=26mm] Q[c,m,wd=26mm] Q[c,m,wd=24mm]}, row{1}={bg=bkVol!14, font=\sffamily\bfseries\footnotesize}, column{1}={font=\sffamily\bfseries\small}, hlines={0.4pt,bkLine}, rowsep=1.6pt, colsep=3pt}
役わり & 力学 & 電気 & 熱 & 水（管とタンク） \\
流れるもの & 速度 $v$ & 電流 $i$ & 熱流 $P$ & 流量 $Q$ \\
押すもの & 力 $F$ & 電圧 $E$ & 温度差 $\Delta T$ & 水位の差 \\
ためるもの & やわらかさ $\dfrac{1}{k}$ & 静電容量 $C$ & 熱容量 $C_{\mathrm{th}}$ & タンク \\
慣性 & 質量 $m$ & インダクタンス $L$ & （ない） & 管の中の水 \\
失うもの & ダンパー $c$ & 抵抗 $R$ & 熱抵抗 $R_{\mathrm{th}}$ & 管の流れにくさ \\
\end{tblr}
\end{bkcenter}
\begin{bkcenter}
\small
\begin{tblr}{colspec={Q[l,m,wd=30mm] Q[c,m,wd=48mm] Q[c,m,wd=48mm]}, row{1}={bg=bkVol!14, font=\sffamily\bfseries\footnotesize}, hlines={0.4pt,bkLine}, rowsep=1.6pt}
 & ばね・質量・ダンパー & RLC 回路 \\
式 & $m\dfrac{\dd v}{\dd t}+cv+kx=F$ & $L\dfrac{\dd i}{\dd t}+Ri+\dfrac{q}{C}=E$ \\
固有角振動数 $\omega_{\mathrm n}$ & $\sqrt{\dfrac{k}{m}}$ & $\dfrac{1}{\sqrt{LC}}$ \\
臨界の値 & $c_{\mathrm c}=2\sqrt{mk}$ & $R_{\mathrm c}=2\sqrt{\dfrac{L}{C}}$ \\
減衰比 $\zeta$ & $\dfrac{c}{2\sqrt{mk}}$ & $\dfrac{R}{2}\sqrt{\dfrac{C}{L}}$ \\
エネルギー & $\dfrac{1}{2}mv^2+\dfrac{1}{2}kx^2$ & $\dfrac{1}{2}Li^2+\dfrac{q^2}{2C}$ \\
\end{tblr}
\end{bkcenter}
{\footnotesize 力 $\leftrightarrow$ 電圧、速度 $\leftrightarrow$ 電流、$m\leftrightarrow L$、$c\leftrightarrow R$、$k\leftrightarrow\dfrac{1}{C}$、$x\leftrightarrow q$。熱には慣性がないので、熱の式は1回微分の式になり、ゆれない（6-15）。}

\bksubsection{8}{電卓・表計算・Python の使い方}
{\small
\begin{longtblr}{colspec={Q[l,m,wd=23mm] Q[l,m,wd=60mm] Q[l,m,wd=31mm] Q[c,m,wd=10mm]}, rowhead=1, row{1}={bg=bkVol!14, font=\sffamily\bfseries\footnotesize}, column{2}={font=\footnotesize}, hlines={0.4pt,bkLine}, rowsep=1.6pt, colsep=3pt}
すること & 書き方（例） & 意味・答え & 章 \\
{オイラー法\\（表計算）} & {C2 \texttt{=-0.05*(B2-20)}\\B3 \texttt{=B2+C2*\$D\$2}\\A3 \texttt{=A2+\$D\$2}} & C列が変わる速さ、D2が $\Delta t$ & 6-2 \\
{速度を先に\\（表計算）} & {D2 \texttt{=-16*B2}\\C3 \texttt{=C2+D2*\$E\$2}\\B3 \texttt{=B2+C3*\$E\$2}} & 列は $t$・$x$・$v$・$a$、E2が $\Delta t$ & {6-2\\6-6} \\
{オイラー法\\（Python）} & \texttt{T = T + (-0.05 * (T - 20)) * dt} & オイラー法の1行 & 6-2 \\
{速度を先に\\（Python）} & {\texttt{v = v + a * dt}\\\texttt{x = x + v * dt}} & この順番で書く & 6-2 \\
$e$ の何乗 & \texttt{=EXP(A2)} & 1なら2.71828… & 6-1 \\
自然対数 & 電卓の ln キー、\ \texttt{=LN(A2)} & 2なら0.693 & {5-14\\6-4} \\
双曲線関数 & {\texttt{=COSH(A2)}\ \ \texttt{=SINH(A2)}\ \ \texttt{=TANH(A2)}\\電卓は hyp を押してから cos など} & 1なら1.543、1.175、0.762 & 6-11 \\
複素数を作る & \texttt{=COMPLEX(3,4,"j")} & \texttt{3+4j} & 6-8 \\
{四則\\（表計算）} & {\texttt{=IMSUM("3+4j","1-2j")}\\\texttt{=IMPRODUCT("3+4j","1-2j")}\\\texttt{=IMDIV("5","1+2j")}} & {\texttt{4+2j}\\\texttt{11-2j}\\\texttt{1-2j}} & 6-8 \\
実部・虚部 & {\texttt{=IMREAL("3+4j")}\\\texttt{=IMAGINARY("3+4j")}} & 3、4 & 6-8 \\
{複素数\\（Python）} & {\texttt{z = 3+4j}\\\texttt{z.real}\ \ \texttt{z.imag}\ \ \texttt{z.conjugate()}} & 3.0、4.0、\texttt{(3-4j)} & 6-8 \\
{平方根\\（Python）} & \texttt{cmath.sqrt(-4)} & \texttt{2j} & 6-8 \\
{絶対値\\（表計算）} & \texttt{=SQRT(B2\textasciicircum 2+C2\textasciicircum 2)} & B が実部、C が虚部 & 6-9 \\
偏角 & {\texttt{=DEGREES(ATAN2(B2,C2))}\\\texttt{=IMARGUMENT("-1-1j")}\\\texttt{cmath.phase(-1-1j)}} & {度で $-135$\\ラジアンで $-2.356$\\ラジアンで $-2.356$} & 6-9 \\
{$e^{j\theta}$\\（Python）} & \texttt{cmath.exp(1j*math.pi/3)} & {\texttt{(0.5+0.866j)}\\（$\frac12+\frac{\sqrt3}{2}j$）} & 6-10 \\
\end{longtblr}\par}
\begin{itemize}[leftmargin=1.4\zw, itemsep=2pt, topsep=4pt, label={\color{\bkTT}●}]
\item $\Delta t$ は1つのセルに入れ、\texttt{\$D\$2} のように \texttt{\$} をつけて、コピーしてもずれないようにします（6-2）。
\item ばねの表で、位置を古い速度 \texttt{C2} で進めると（ふつうのオイラー法）、ゆれがふくらみます。新しい速度 \texttt{C3} を使います（6-2）。
\item Python では、$j$ そのものは \texttt{1j} と書きます。\texttt{j} だけだと文字の名前になります（6-8）。
\item \texttt{ATAN2} は横（実部）を先に入れます（6-9）。答えの角はラジアンなので、度にするには \texttt{DEGREES} をかぶせます。
\item \texttt{cmath.exp(1j*math.pi)} が \texttt{(-1+1.22e-16j)} になるのは、コンピュータの丸めのずれです。ほんとうは $-1$ です（6-10）。
\end{itemize}

\Needspace{8\baselineskip}
\bkbkhead{プログラムのひな形（6-2）}
{\small 左は冷めるお湯（$\Delta t=0.1$ 分で10分後。4行めがオイラー法の1行。式の答えは56.39\,℃）。右はばね（速度を先に進めるオイラー法。$\Delta t=0.02$\,s で5秒後の位置 \texttt{0.03665…} を表示。式の答えは $0.1\cos(4\times5)\approx0.0408$）。右の \texttt{v =} と \texttt{x =} の行を入れかえると、ふつうのオイラー法になり、ゆれがふくらみます。}\par\vspace{1mm}
\noindent\begin{minipage}[t]{0.485\linewidth}\vspace{0pt}
\begin{bkcode}[Python]
T = 80.0
dt = 0.1
for i in range(100):
    T = T + (-0.05*(T-20))*dt
print(T)
\end{bkcode}
\begin{bkcodeout}
56.34622618944369
\end{bkcodeout}
\end{minipage}\hfill
\begin{minipage}[t]{0.485\linewidth}\vspace{0pt}
\begin{bkcode}[Python]
x = 0.1
v = 0.0
dt = 0.02
for i in range(1, 251):
    a = -16 * x
    v = v + a * dt
    x = x + v * dt
print(x)
\end{bkcode}
\end{minipage}\par\vspace{2mm}
%% ---------------- チェック表 ----------------
\bkbacktitle{できるようになったことチェック表}
{\small 各章の「ふりかえり」を集めた表です。章を学び終えたら1回め、次の巻に進む前に2回めのチェックをしよう。できなかった所は、ページにもどって読みなおします。}\par\vspace{2mm}
%% 自動生成（tools/mkchecklist.py）。各章の「ふりかえり」の項目を集めたもの
\begingroup\footnotesize
\def\bkcb{\tikz[baseline=-0.5ex]\draw[\bkTF, line width=0.7pt, rounded corners=0.5pt] (0,-0.3em) rectangle (0.8em,0.5em);}
\begin{longtblr}{colspec={Q[c,m,wd=9mm] X[l,m] Q[c,m,wd=10mm] Q[c,m,wd=9mm] Q[c,m,wd=9mm]}, width=\linewidth, rowhead=1, colsep=3pt,
  row{1}={bg=bkVol!14, font=\sffamily\bfseries\scriptsize}, hlines={0.4pt,bkLine}, rowsep=0.7pt}
章 & できるようになったこと & ページ & 1回め & 2回め \\
{\sffamily\bfseries\color{\bkTT}6-1} & 微分方程式を、ことばと単位で読み、式に書ける。 & \pageref{bk:6-1:yomu} & \bkcb & \bkcb \\
 & 傾きの場から、解のようすを予想できる。 & \pageref{bk:6-1:ba} & \bkcb & \bkcb \\
 & 代入して、解かどうか確かめられる。 & \pageref{bk:6-1:tashikame} & \bkcb & \bkcb \\
{\sffamily\bfseries\color{\bkTT}6-2} & 「次の値＝今の値＋今の変わる速さ×$\Delta t$」で、表計算で解ける。 & \pageref{bk:6-2:ippo} & \bkcb & \bkcb \\
 & $\Delta t$ を小さくすると正確になることを確かめられる。 & \pageref{bk:6-2:zure} & \bkcb & \bkcb \\
 & ばねの式を、位置と速度の2つの列で、正しい順に解ける。 & \pageref{bk:6-2:bane} & \bkcb & \bkcb \\
{\sffamily\bfseries\color{\bkTT}6-3} & 右辺を見て、積分するだけで解ける形を見分け、解ける。 & \pageref{bk:6-3:miwake} & \bkcb & \bkcb \\
 & 深さや場所が変数のルールも、積分して解ける。 & \pageref{bk:6-3:basho} & \bkcb & \bkcb \\
 & 両はしの条件（境界条件）から、積分定数を決められる。 & \pageref{bk:6-3:kyokai} & \bkcb & \bkcb \\
{\sffamily\bfseries\color{\bkTT}6-4} & 変数分離で、自分に比例してへる式を解ける。 & \pageref{bk:6-4:bunri} & \bkcb & \bkcb \\
 & 時定数・半減期を使って、かかる時間を求められる。 & \pageref{bk:6-4:tau} & \bkcb & \bkcb \\
 & RC・RL の回路の式を立てて、時定数を求められる。 & \pageref{bk:6-4:rc} & \bkcb & \bkcb \\
{\sffamily\bfseries\color{\bkTT}6-5} & 落ちる物の式を立て、終端速度を求められる。 & \pageref{bk:6-5:shuutan} & \bkcb & \bkcb \\
 & 表計算で、落ちる物の速さを計算できる。 & \pageref{bk:6-5:hyou} & \bkcb & \bkcb \\
 & 速さに比例する抵抗の式を解ける（RL 回路も）。 & \pageref{bk:6-5:toku} & \bkcb & \bkcb \\
{\sffamily\bfseries\color{\bkTT}6-6} & ばねの運動方程式を立て、$\omega_{\mathrm n}=\sqrt{k/m}$ と周期を求められる。 & \pageref{bk:6-6:shiki} & \bkcb & \bkcb \\
 & はじめの位置と速度から、振幅と初期位相を決められる。 & \pageref{bk:6-6:shoki} & \bkcb & \bkcb \\
 & 振り子・LC回路の周期を、ばねと同じ形の式から求められる。 & \pageref{bk:6-6:lc} & \bkcb & \bkcb \\
{\sffamily\bfseries\color{\bkTT}6-7} & ばね・質量・ダンパーの式を、力の図から立てられる。 & \pageref{bk:6-7:shiki} & \bkcb & \bkcb \\
 & となりの山の比から、対数減衰率を求められる。 & \pageref{bk:6-7:yama} & \bkcb & \bkcb \\
 & 減衰比 $\zeta$ を計算して、止まり方を答えられる（RLC回路でも）。 & \pageref{bk:6-7:zeta} & \bkcb & \bkcb \\
{\sffamily\bfseries\color{\bkTT}6-8} & $j$ をかけると90°回ることから、$j^2=-1$ を説明できる。 & \pageref{bk:6-8:kaiten} & \bkcb & \bkcb \\
 & 複素数のたし算・ひき算・かけ算・わり算ができる。 & \pageref{bk:6-8:kakeru} & \bkcb & \bkcb \\
 & 判別式がマイナスの2次方程式を解いて、虚数の解を求められる。 & \pageref{bk:6-8:kai} & \bkcb & \bkcb \\
{\sffamily\bfseries\color{\bkTT}6-9} & 複素数を平面の点で表し、絶対値と偏角を求められる。 & \pageref{bk:6-9:kyoku} & \bkcb & \bkcb \\
 & かけ算が「大きさはかけ算、角はたし算」になることを説明できる。 & \pageref{bk:6-9:kakeru} & \bkcb & \bkcb \\
 & かけ算で、点を回したり、うでの先の位置を求めたりできる。 & \pageref{bk:6-9:kaiten} & \bkcb & \bkcb \\
{\sffamily\bfseries\color{\bkTT}6-10} & $e^{j\theta}$ が単位円の上の角 $\theta$ の点になるわけを、小さな回転のくり返しと微分で説明できる。 & \pageref{bk:6-10:kaiten} & \bkcb & \bkcb \\
 & $e^{j\theta}$ を $a+bj$ に直し、$re^{j\theta}$ のかけ算・わり算・べきを計算できる。 & \pageref{bk:6-10:kyoku} & \bkcb & \bkcb \\
 & フェーザのたし算で交流を合成し、微分が $j\omega$ をかけることだと説明できる。 & \pageref{bk:6-10:phasor} & \bkcb & \bkcb \\
{\sffamily\bfseries\color{\bkTT}6-11} & $\cosh x$・$\sinh x$・$\tanh x$ を $e^x$ と $e^{-x}$ で表し、値とグラフの形を説明できる。 & \pageref{bk:6-11:teigi} & \bkcb & \bkcb \\
 & $\cosh^2x-\sinh^2x=1$ と、微分のきまりを使える。 & \pageref{bk:6-11:bibun} & \bkcb & \bkcb \\
 & $y''=a^2y$ の答えを $\cosh$・$\sinh$ で書き、落下や電線の形に使える。 & \pageref{bk:6-11:kaku} & \bkcb & \bkcb \\
{\sffamily\bfseries\color{\bkTT}6-12} & 式に $e^{st}$ を入れて、特性方程式を作れる。 & \pageref{bk:6-12:oku} & \bkcb & \bkcb \\
 & 解が2つ・複素数・重解のとき、答えを組み立てられる。 & \pageref{bk:6-12:ba3} & \bkcb & \bkcb \\
 & $\zeta$ で動き方を見分け、車や RLC 回路の数を求められる。 & \pageref{bk:6-12:zeta} & \bkcb & \bkcb \\
{\sffamily\bfseries\color{\bkTT}6-13} & ゆすり続けたときのゆれを、表計算と共振曲線で説明できる。 & \pageref{bk:6-13:kyoshin} & \bkcb & \bkcb \\
 & ゆすりに合うゆれを、係数を比べる方法と複素数で求められる。 & \pageref{bk:6-13:keisu} & \bkcb & \bkcb \\
 & いちばん大きくゆれる $\omega$ を微分で求め、RLC 回路の共振周波数を求められる。 & \pageref{bk:6-13:yama} & \bkcb & \bkcb \\
{\sffamily\bfseries\color{\bkTT}6-14} & からみあう2つの量を、表計算で同時に進めて計算できる。 & \pageref{bk:6-14:hyou} & \bkcb & \bkcb \\
 & 和と差（重みつきの和と差）に分けて、2つの式を解ける。 & \pageref{bk:6-14:wasa} & \bkcb & \bkcb \\
 & 振動モードや消去法で、連立微分方程式の答えを求められる。 & \pageref{bk:6-14:mode} & \bkcb & \bkcb \\
{\sffamily\bfseries\color{\bkTT}6-15} & 法則から、微分方程式とはじめの状態を書ける。 & \pageref{bk:6-15:buhin} & \bkcb & \bkcb \\
 & 相似の表で、ばね・電気・熱の量を対応させられる。 & \pageref{bk:6-15:soji} & \bkcb & \bkcb \\
 & 式の形から、解き方（章）を選べる。 & \pageref{bk:6-15:chizu} & \bkcb & \bkcb \\
\end{longtblr}
\endgroup


%% ---------------- つぎの巻へ（第7巻の予告と、この巻の道具がこの先どこで使われるか） ----------------
\bkbacktitle{つぎの巻へ}
\newcommand\bknextpill[1]{\bk@pill{#1}}
\def\bknextrow{1}
\newcommand\bknextcard[4]{%
  \begin{tcolorbox}[enhanced, nobeforeafter, box align=top, width=0.485\linewidth, colframe=bkVol7, colback=white, boxrule=0.9pt, arc=2.5mm,
    left=2.5mm, right=2.5mm, top=1.5mm, bottom=2mm, equal height group=bknext\bknextrow]
  \bknextpill{#1}\enspace{\sffamily\bfseries\small #2}\par\vspace{1mm}
  {\sffamily\bfseries\footnotesize\color{bkDeep7}#3}\par\vspace{1mm}
  {\footnotesize\setlength{\parindent}{0pt}{\sffamily\bfseries\color{bkInk}使う道具}\enspace #4}
  \end{tcolorbox}}
未来を計算する道具がそろいました。第7巻「まとめて計算する」では、たくさんの数を表（行列）のまま計算します。6-14 で並べた係数の表が、その入口です。第7巻の章をのぞいてみましょう。\par\vspace{2.5mm}
\noindent\bknextcard{7-1}{数の表で計算する：行列}{買い物の合計を、表のかけ算で一度に求める}{2つの式の係数を、表の形に並べたこと（6-14）。}\hfill
\bknextcard{7-6}{回転行列：くるっと回す}{30°回して45°回すと、75°回したのと同じ}{複素数のかけ算は回転（6-9）。}\par\vspace{1.2mm}
\def\bknextrow{2}\noindent\bknextcard{7-8}{固有値と固有ベクトル：向きの変わらない方向}{変換しても、向きが変わらない矢印がある}{特性方程式（6-12）と、振動モード（6-14）。}\hfill
\bknextcard{7-9}{対角化：まっすぐな方向に分けて考える}{こみいった変換も、のびるだけの変換に分けられる}{和と差に分けると、連立の式がほどけること（6-14）。}\par\vspace{1.2mm}
\def\bknextrow{3}\noindent\bknextcard{7-10}{行列のくり返し：漸化式とオイラー法を行列で}{何回もかけるとどうなるかは、固有値が教えてくれる}{オイラー法（6-2）と、連立の表計算（6-14）。}\hfill
\bknextcard{7-11}{行列の使いみち}{ロボット、電気回路、画像…行列はどこで働いている？}{ばね・電気・熱の相似の表（6-15）。}\par
\bkfukidashi{delta}{第6巻の道具は、第7巻・第9巻・物理シリーズでずっと使うよ。自信のない章を見なおしてから、第7巻へ進もう！}

\clearpage
\bkbkhead{この巻の道具は、この先でこう使う}
{\small 第6巻の章と、その道具を使う第7巻から先の章です（章のはじめの「道しるべ」の「この章の先に」から集めました）。}\par\vspace{3mm}
%% 自動生成（tools/mknext.py）。各章の道しるべ「この章の先に」から、後ろの巻の章を集めたもの
\begin{longtblr}{colspec={Q[l,m,wd=58mm] X[l,m]}, width=\linewidth, rowhead=1, row{1}={bg=bkVol!14, font=\sffamily\bfseries\footnotesize},
  hlines={0.4pt,bkLine}, rowsep=1.5pt, cells={font=\footnotesize}}
この巻の章 & この道具を使う章 \\
{\bknextpill{6-2}\ 表計算で解こう：オイラー法} & {\bknextpill{7-10}\ 行列のくり返し：漸化式とオイラー法を行列で} \\
{\bknextpill{6-3}\ 積分するだけで解ける形} & {\bknextpill{8-13}\ 熱と波の伝わり方：偏微分方程式} \\
{\bknextpill{6-4}\ 自分に比例してへる：冷めるお湯} & {\bknextpill{9-5}\ 微分をかけ算に変える翻訳機：ラプラス変換\\ \bknextpill{9-9}\ 極と安定性：s の値でふるまいがわかる} \\
{\bknextpill{6-6}\ ばねと振り子} & {\bknextpill{7-8}\ 固有値と固有ベクトル：向きの変わらない方向} \\
{\bknextpill{6-8}\ 2乗してマイナスになる数：虚数} & {\bknextpill{9-9}\ 極と安定性：s の値でふるまいがわかる} \\
{\bknextpill{6-9}\ 複素平面：かけ算は回転と拡大} & {\bknextpill{7-6}\ 回転行列：くるっと回す} \\
{\bknextpill{6-10}\ オイラーの公式} & {\bknextpill{9-1}\ 波をたし合わせて形を作る：フーリエ級数\\ \bknextpill{9-3}\ 複素数で書くフーリエ級数} \\
{\bknextpill{6-11}\ 双曲線関数 cosh・sinh} & {\bknextpill{8-13}\ 熱と波の伝わり方：偏微分方程式} \\
{\bknextpill{6-12}\ 減衰振動を式で解く：特性方程式} & {\bknextpill{9-7}\ ラプラス変換で微分方程式を解く\\ \bknextpill{9-9}\ 極と安定性：s の値でふるまいがわかる} \\
{\bknextpill{6-13}\ ゆらし続けると：強制振動と共振} & {\bknextpill{9-6}\ 変換表と性質\\ \bknextpill{9-7}\ ラプラス変換で微分方程式を解く\\ \bknextpill{9-8}\ 伝達関数：入れると出てくるを s で表す\\ \bknextpill{9-12}\ フーリエとラプラスの使いみち} \\
{\bknextpill{6-14}\ 微分方程式がいくつも：連立微分方程式} & {\bknextpill{7-8}\ 固有値と固有ベクトル：向きの変わらない方向\\ \bknextpill{7-9}\ 対角化：まっすぐな方向に分けて考える\\ \bknextpill{7-10}\ 行列のくり返し：漸化式とオイラー法を行列で\\ \bknextpill{9-10}\ 連立微分方程式を行列で：状態方程式\\ \bknextpill{9-11}\ 行列の指数関数と固有値} \\
{\bknextpill{6-15}\ まとめ：物理は「微分方程式＋はじめの状態」で書ける} & {\bknextpill{7-0}\ 表のままで計算できたら？\\ \bknextpill{8-13}\ 熱と波の伝わり方：偏微分方程式\\ \bknextpill{9-5}\ 微分をかけ算に変える翻訳機：ラプラス変換\\ \bknextpill{9-10}\ 連立微分方程式を行列で：状態方程式} \\
\end{longtblr}


%% ---------------- さくいん ----------------
\bkbacktitle{さくいん}
{\small 数字はページです。その語がはじめて出てくるページや、くわしく説明しているページをのせています。}\par\vspace{2mm}
\begin{multicols}{2}\small\setlength{\parskip}{0pt}
\directlua{bk.printidx()}
\end{multicols}

%% ---------------- 奥付 ----------------
\bkpadfour
\global\bk@decofalse\bkfolio{false}
\null\vfill
\noindent\begin{minipage}{\linewidth}\small
{\sffamily\bfseries\Large\color{bkDeep}物理を理解するための数学　第6巻}\par\vspace{1mm}
{\sffamily\bfseries\large 未来を計算する}\par\vspace{3mm}
{\color{bkVol}\hrule height 1pt}\vspace{3mm}
\begin{tabular}{@{}l@{\hspace{6mm}}l@{}}
編者 & 「物理を理解するための数学」編集委員会 \\
 & 委員長　小学校長（まとめ役） \\
 & 委員　数学・電気工学・機械工学・教育学・量子力学・建築工学・\\
 & \phantom{委員　}半導体工学の大学教授、幼稚園教諭、塾講師、\\
 & \phantom{委員　}特別支援学校教諭、心理士、科学解説者、小学校教諭 \\
デザイン & ユニバーサルデザイン担当委員 \\
挿絵 & イラスト担当委員 \\
版 & 第1版（第14回編集委員会）　2026年10月 \\
\end{tabular}\par\vspace{3mm}
{\color{bkVol}\hrule height 1pt}\vspace{2mm}
{\footnotesize この本の問題の答えは、巻末の「こたえ」にあります。くわしい解き方は、別冊の問題集と指導書にのせます。}
\end{minipage}
\makeatother
